% Generated draft: requires mathematical and source review.
\documentclass[11pt]{article}
\usepackage[margin=25mm]{geometry}
\usepackage{fontspec}
\setmainfont{Latin Modern Roman}
\usepackage{amsmath,amssymb,mathtools,unicode-math}
\setmathfont{Latin Modern Math}
\usepackage{enumitem,needspace,xurl,hyperref}
\setlength{\parindent}{0pt}
\setlength{\parskip}{5pt}
\setlength{\emergencystretch}{3em}
\newcommand{\N}{\mathbb{N}}
\newcommand{\Z}{\mathbb{Z}}
\newcommand{\Q}{\mathbb{Q}}
\newcommand{\R}{\mathbb{R}}
\newcommand{\C}{\mathbb{C}}
\newcommand{\F}{\mathbb{F}}
\newcommand{\A}{\mathbb{A}}
\newcommand{\PP}{\mathbb{P}}
\newcommand{\E}{\mathbb{E}}
\newcommand{\eps}{\varepsilon}
\newcommand{\dd}{\,\mathrm{d}}
\begin{document}
\section*{UnsolvedMath ID: 30002862}
\textbf{Source problem code:} catalogue entry
\subsection{Definitions and mathematical statement}
The supplied record leaves the dimension \(n\) and the matrix \(M\) unspecified. To interpret its displayed expressions, one needs a fixed finite dimension \(n\) and a fixed real matrix \(M\in\R^{n\times n}\). These are typing requirements; the record supplies no construction of \(M\), hypotheses on its entries, or intended class of matrices.

For vectors \(u,v\in\R^n\), write \(u\geq0\) when \(u_i\geq0\) for every coordinate, and define the Euclidean inner product by \(\langle u,v\rangle=\sum_{i=1}^n u_iv_i\). Conditional on the intended fixed choices of \(n\) and \(M\), the supplied question is whether
\[
\forall q\in\R^n\;\exists x\in\R^n:\qquad
x\geq0,\qquad Mx+q\geq0,\qquad
\langle Mx+q,x\rangle=0.
\]
Equivalently, the last condition, together with the two coordinatewise inequalities, requires \(x_i(Mx+q)_i=0\) for every coordinate \(i\).

The record calls this the complementarity form of the Schumaker--Strang conjecture, but supplies neither the matrix arising from that conjecture nor a statement establishing the correspondence. Thus the formula records the stated complementarity question, while the intended mathematical conjecture remains underdetermined. In particular, no universal quantifier over matrices \(M\) is supplied.
\subsection{Short English statement}
For an intended fixed matrix \(M\), the question asks whether every vector \(q\) admits a nonnegative vector \(x\) such that \(Mx+q\) is nonnegative and each coordinate of \(x\) has zero product with the corresponding coordinate of \(Mx+q\). The matrix and its hypotheses are missing, so the intended conjecture cannot yet be stated precisely.
\subsection{Sources}
\begin{itemize}
\item Local UnsolvedMath registry, supplied record with permanent ID \(30002862\), source problem code \texttt{catalogue entry}, titled ``Complementarity Form of the Schumaker Strang Conjecture.'' The record belongs to the Oberwolfach Reports Open Problems collection and gives proposed year 2015; no individual proposer is identified.
\item Supplied report citation: \emph{Multivariate Splines and Algebraic Geometry} (2016), Oberwolfach Reports, DOI 10.4171/owr/2015/21, pp. 1139--1200. \url{https://doi.org/10.4171/owr/2015/21}. The background describes Oberwolfach Reports as documenting mathematical workshops and their open-problem discussions. The supplied literature-source label for this link is ``Oberwolfach Report 21.''
\item Supplied literature source: ``Generalized Splines on Graphs, EJC 2024.'' \url{https://www.combinatorics.org/ojs/index.php/eljc/article/download/v31i1p29/pdf}.
\item Literature assessment dated exactly \texttt{2026-08-21}: the supplied curation labels the problem open and states that the Schumaker--Strang lower-bound conjecture and its complementarity formulation remain active in general, with important degree and genericity cases proved but without a proof of the stated full conjecture. This is supplied dated curation derived from the OWR material, not an independently verified literature assessment. Its time-sensitive status and cited claims require independent verification.
\end{itemize}
\end{document}
